\documentclass[11pt]{article}
\usepackage[a4paper,margin=27mm]{geometry}
\usepackage{amsmath,amssymb,amsthm,enumitem}
\usepackage[hidelinks]{hyperref}
\newtheorem{theorem}{Theorem}
\newtheorem{lemma}{Lemma}
\newcommand{\R}{\mathbb R}
\title{Order-independent Hermite means:\\a counterexample and the corrected classification}
\author{Principia Math In Progress}
\date{12 September 2026}
\begin{document}
\maketitle
\begin{abstract}
Conjecture 23 of Horwitz is false: $f(x)=1/(x+1)$ produces the same
mean for every admissible pair of interpolation orders, but that mean
is not the harmonic mean. We also classify all generators satisfying
the source's nondegeneracy hypotheses at every order: they are shifted
reciprocals plus polynomials of degree at most two.
\end{abstract}

\section*{Task 1. The problem and its full resolution}
For $0<a<b$ and integers $0\le m_2<m_1$, put
$n=m_1+m_2+1$. Let $P,Q$ be the degree-at-most-$n$ Hermite
interpolants to $f$: $P$ matches derivatives through orders $m_1,m_2$
at $a,b$, while $Q$ interchanges these orders. Write $m=m_1-m_2$.
The source defines $M_{f,m_1,m_2}(a,b)$ as the unique interior solution of
$P=Q$ for even $m$, and of $2f=P+Q$ for odd $m$.
Equivalently, with $E_P=f-P$ and $E_Q=f-Q$, the equation is
\begin{equation}\label{eq:errors}
 E_P(x)=(-1)^m E_Q(x).
\end{equation}
The existence hypotheses are $f\in C^{n+2}$ and
$f^{(n+1)}\ne0$ on every relevant positive interval.
Conjecture 23, on published page 95 of \cite{horwitz}, asserts that
independence of both orders characterizes $f(x)=C/x$.

\begin{theorem}\label{thm:classification}
Suppose $f\in C^\infty(0,\infty)$ and
$f^{(r)}(x)\ne0$ for every $r\ge3$ and $x>0$.
Its Hermite means are independent of $m_1,m_2$ if and only if
\begin{equation}\label{eq:family}
 f(x)=\frac{C}{x+c}+p(x),\qquad
 C\in\R\setminus\{0\},\quad c\ge0,\quad \deg p\le2.
\end{equation}
Their common value is
\begin{equation}\label{eq:mean}
 M_c(a,b)=\frac{2ab+c(a+b)}{a+b+2c}
 =\frac{2(a+c)(b+c)}{a+b+2c}-c.
\end{equation}
In particular, $f(x)=1/(x+1)$ is a counterexample to Conjecture 23.
\end{theorem}

\subsection*{Exact proof of sufficiency and of the counterexample}
First take $f(x)=C/(x+c)$ and set
\[
 W_P(x)=(x-a)^{m_1+1}(x-b)^{m_2+1},\qquad
 W_Q(x)=(x-a)^{m_2+1}(x-b)^{m_1+1}.
\]
The polynomial $C-(x+c)P(x)$ has degree at most $n+1$ and vanishes
at the nodes with exactly the required multiplicities, or higher ones.
It is therefore a constant multiple of the monic polynomial $W_P$.
Evaluation at $x=-c$ determines that constant and gives
\begin{equation}\label{eq:exact-remainder}
 E_P(x)=\frac{C W_P(x)}{(x+c)W_P(-c)},\qquad
 E_Q(x)=\frac{C W_Q(x)}{(x+c)W_Q(-c)}.
\end{equation}
The denominators are nonzero since $a,b>0$ and $c\ge0$.
For $a<x<b$, their ratio is
\[
 \frac{E_P(x)}{E_Q(x)}
 =(-1)^m\left(\frac{(x-a)(b+c)}{(b-x)(a+c)}\right)^m.
\]
The quantity in parentheses is positive. Thus \eqref{eq:errors}, for
either parity and every $m\ge1$, is equivalent to
\[
 \frac{x-a}{b-x}=\frac{a+c}{b+c}.
\]
This has exactly the solution \eqref{eq:mean} in $(a,b)$.
Every allowed degree $n$ is at least two, so adding a fixed quadratic
$p$ adds $p$ to both interpolants and leaves both errors unchanged.
This proves sufficiency. The family has
$f^{(r)}(x)=(-1)^r r!C/(x+c)^{r+1}\ne0$ for $r\ge3$.

For $c=1,C=1,p=0$, the function is not $C_0/x$ for any constant
$C_0$, since $x/(x+1)$ is not constant. Its mean differs from the
harmonic mean $H$ and the arithmetic mean $A$ whenever $a\ne b$:
\begin{equation}\label{eq:differences}
 M_c-H=\frac{c(a-b)^2}{(a+b)(a+b+2c)},\qquad
 A-M_c=\frac{(a-b)^2}{2(a+b+2c)}.
\end{equation}
For example, $M_1(1,2)=7/5$, whereas $H(1,2)=4/3$.
Thus the counterexample changes the common mean itself. The source
already notes invariance under polynomial additions in Remark 5;
that invariance alone is not the basis of this counterexample.

\subsection*{A local coefficient of an arbitrary Hermite mean}
For confluent divided differences, let $a^{[j]}$ denote $j$ copies
of the node $a$. The Hermite error formula turns \eqref{eq:errors} into
\begin{equation}\label{eq:divided}
 (u-a)^m D_P=(b-u)^m D_Q,
\quad D_P=f[u,a^{[m_1+1]},b^{[m_2+1]}],
\quad D_Q=f[u,a^{[m_2+1]},b^{[m_1+1]}].
\end{equation}
This is equation (4) of the source and is valid for both parities.

\begin{lemma}\label{lem:diagonal}
Fix an allowed pair and $t>0$. For smooth $f$ with
$f^{(n+1)}(t)\ne0$,
\[
 M_{f,m_1,m_2}(t-h,t+h)
 =t+\frac{f^{(n+2)}(t)}{(n+2)f^{(n+1)}(t)}h^2+O(h^3)
 \qquad(h\downarrow0).
\]
\end{lemma}
\begin{proof}
The divided difference of order $N$ at nodes $t+\delta_0,\ldots,t+\delta_N$
has the expansion
\begin{equation}\label{eq:dd-expansion}
 \frac{f^{(N)}(t)}{N!}
 +\frac{f^{(N+1)}(t)}{(N+1)!}\sum_{i=0}^N\delta_i+O(h^2)
 \quad\text{if }|\delta_i|\le h.
\end{equation}
For completeness, the Hermite--Genocchi integral writes this divided
difference as $\mathbb E f^{(N)}(t+\sum_i V_i\delta_i)/N!$, where
$V$ is uniform on the simplex $\sum_iV_i=1$, $V_i\ge0$.
Taylor expansion and $\mathbb E V_i=1/(N+1)$ give
\eqref{eq:dd-expansion}, including repeated nodes, with uniform remainder.
The integral itself follows by repeated use of the fundamental theorem
of calculus.

Write $u=t+s$, so $|s|<h$, and put
$B=f^{(n+1)}(t)/(n+1)!\ne0$ and
$D=f^{(n+2)}(t)/(n+2)!$.
The sums of offsets in the two divided differences are $s-mh$ and $s+mh$.
Consequently
\[
 D_P=B+D(s-mh)+O(h^2),\qquad
 D_Q=B+D(s+mh)+O(h^2).
\]
For small $h$, their ratio is positive and equals $1+O(h)$.
Equation \eqref{eq:divided} first implies
$m\log((h+s)/(h-s))=O(h)$, hence $s=O(h^2)$.
Expanding once more gives
\[
 2m\frac{s}{h}+O(h^3)
 =\log\frac{D_Q}{D_P}
 =2mh\frac{D}{B}+O(h^2).
\]
Multiplication by $h/(2m)$ proves the formula.
\end{proof}

\subsection*{Necessity and the global domain}
If the means are independent of the orders, compare $(m_1,m_2)=(1,0)$
and $(2,0)$, whose degrees are $n=2$ and $n=3$.
Lemma~\ref{lem:diagonal} implies, at every $t>0$,
\[
 \frac{f^{(4)}}{4f^{(3)}}=\frac{f^{(5)}}{5f^{(4)}}.
\]
Set $v=f^{(4)}/(4f^{(3)})$, which is smooth and never zero.
Differentiating and using the displayed identity gives
\[
 v'=\frac{f^{(5)}}{4f^{(3)}}
       -\frac{(f^{(4)})^2}{4(f^{(3)})^2}
    =5v^2-4v^2=v^2.
\]
Thus $(-1/v)'=1$, so $v(t)=-1/(t+c)$ for a constant $c$.
The denominator cannot vanish anywhere on $(0,\infty)$; hence $c\ge0$.
The equation $f^{(4)}/f^{(3)}=-4/(t+c)$ now yields
$f^{(3)}=K/(t+c)^4$ with $K\ne0$. Integrating three times gives
\eqref{eq:family}, with $C=-K/6$. This proves the theorem.

The nondegeneracy in the classification is precisely what results from
requiring the source's hypotheses for every allowed pair: $n+1$ runs
through all integers at least three. The counterexample satisfies these
hypotheses on the entire positive half-line. Homogeneity is defined as
an additional property in the source, not a hypothesis of Conjecture 23.
If homogeneity is imposed additionally, \eqref{eq:mean} forces $c=0$.

\section*{Follow-up tasks}
These tasks follow the full resolution of Task 1 and are not additional
completed or independently screened open problems.
\paragraph{Task 2.} Classify order-independent generators on arbitrary real intervals,
including intervals on either side of a pole and the appropriate
conditions at finite endpoints.
\paragraph{Task 3.} Determine which two prescribed order pairs already force a shifted
reciprocal plus a polynomial, and identify exceptional pairs.
\paragraph{Task 4.} Prove quantitative stability estimates: bound a generator's distance
from the classified family using a uniform bound on the difference
between its means for two consecutive total degrees.
\paragraph{Task 5.} Classify generators whose means become independent of the orders
only after a fixed total degree, and determine the resulting polynomial
freedom and possible limiting means.

\section*{Status and verification}
The dated accompanying files record the live open status, the search for
earlier public solution claims, and exact checks of confluent interpolation,
both parity equations, and the local coefficient. The proof is analytic;
the finite checks are supplementary. No independent peer review or
proof-assistant certification is claimed.

\begin{thebibliography}{9}
\bibitem{horwitz} A. Horwitz, \emph{Means and Hermite Interpolation},
Journal of Mathematical Inequalities \textbf{2}(1) (2008), 75--95,
Conjecture 23. \url{https://doi.org/10.7153/jmi-02-08}.
\url{https://files.ele-math.com/articles/jmi-02-08.pdf}.
Preprint: \url{https://arxiv.org/abs/0711.4940}.
\end{thebibliography}
\end{document}
